% !PRINTED 409
\weberpage{409}
For the case of an \gesperrt{even} $\Delta$ one decomposes $k$ into

\[
v, \quad -\frac{\Delta}{2} + v, \qquad 0 < v < \frac{\Delta}{2}
\]

and applies the formula:

\[
\left(\Delta,\ -\frac{\Delta}{2} + v\right) = -(\Delta,v),
\]

which follows easily from the theorems of § 85.

Then one finds first:

\[
h = \frac{1}{2}\overset{v}{\Sigma}(\Delta,v),
\]

and if one decomposes the $v$ once more into

\weq{12}{\begin{aligned}
&\mu, \quad -\frac{\Delta}{2} - \mu: \qquad 0 < \mu \leqq \frac{-\Delta}{4},\\
&h = \Sigma(\Delta,\mu).
\end{aligned}}

The formula for positive discriminants can be transformed by means
of the theory of the division of the circle, which we do not wish to go
into here.

\begin{center}
§ 113.\quad \gesperrt{The Number of Genera.}
\end{center}

Formula (5) of the preceding section permits the determination
 of the exact number of genera after \gesperrt{Dirichlet}.
By $\delta$ we understand any stem divisor of $D$, and we set:

\weq{1}{\begin{aligned}
&F(z) = \frac{(\delta,z)}{z^s}, \quad \text{where } z \text{ is relatively prime to } D\text{,}\\
&F(z) = 0, \qquad \text{if } z \text{ and } D \text{ have a common divisor.}
\end{aligned}}

Then $(\delta,N\mathfrak{a})$ is one of the characters of the genus determined
by $\mathfrak{a}$. If we therefore denote this character by $(\delta,A)$,
it follows from § 112, (5):

\weq{2}{\overset{A}{\Sigma}(\delta,A)\overset{\mathfrak{a}}{\Sigma}\frac{1}{(N\mathfrak{a})^s} = \Sigma\frac{(\delta,m)}{m^s}\Sigma\frac{(\delta',n)}{n^s}.}

In this, $\delta' = \delta\Delta$, and since $n$ is relatively prime to $D$, it can
also be replaced by the \gesperrt{stem} of $\delta\Delta$ (according to the symbolic
 multiplication in § 104). If $\delta = 1$, then $\delta' = \Delta$,
and if $\delta = \Delta$, then $\delta' = 1$.

Apart from the two cases $\delta = 1$, $\delta = \Delta$, the sums on the
right-hand side of (2) therefore retain a finite value for $s = 1$
(which, according to the preceding section, are expressible by class numbers
 of the discriminants $\delta$, $\delta'$).

On the left-hand side of (2), in the sum over $\mathfrak{a}$
the symbol $\mathfrak{a}$ runs through the totality of the ideals of the individual \gesperrt{classes},
%CONT>>
